\documentclass[11pt,a4paper]{article}
\usepackage[margin=24mm]{geometry}
\usepackage{amsmath,amssymb,booktabs,microtype}
\usepackage[hidelinks]{hyperref}
\newcommand{\dd}{\mathrm d}
\title{Route G2-M: A Frozen Readout Contract}
\author{Momentum first; a noisy pump tag only as a conditional alternative}
\date{23 September 2026}
\begin{document}
\maketitle
\begin{abstract}
The selected design uses the existing driven model with incoming width
$\sigma_i=.1$ and momentum resolution $\sigma_d=.05$, without requiring
a sideband tag. We freeze its decision rule before testing calibration
errors. A separate noisy-tag channel quantifies what additional
information would be required to retain coarser momentum resolution.
This is a readout requirement in model units, not a built apparatus.
\end{abstract}

\section{The physical choice and its boundary}
All dynamics, free masses, $g=.05$, $\epsilon=.5$, $\Omega=.2$,
finite internal preparation and the classical truncated-normal beam
are inherited unchanged from G2-S/D. Select the same counted collisions
with $m=0$, $r=|j|=1$. Write $s=\operatorname{sign}j$ and
$l_+=0$, $l_-=-2$. Their unnormalized component densities are
\begin{equation}
 h_{sq}(x)=w_{l_s}\int_0^\infty\dd p\,f_{.1}(p)
 K_{l_s,0,q}(p)\,
 {\cal N}\bigl(x;p_{f,sq}(p),\sigma_d^2\bigr),\qquad
 h_s(x)=\sum_qh_{sq}(x).
\end{equation}
Let $Z=\sum_s\int h_s\,\dd x$, $\pi_s=Z^{-1}\int h_s\,\dd x$.
The same priors must be used throughout the calibration stress test:
instrument response changes the readout, not the preselected event rate.
Perfect mass labels, common COM frame, equal acceptance and no
background contamination remain explicit assumptions.

The main choice is \emph{momentum-only readout}, nominal
$\sigma_d=.05$. It needs fewer new physical assumptions than an
independent event-by-event measurement of pump work. Its nominal rule is
\begin{equation}\label{eq:frozen}
 \delta_0(x)=
 \begin{cases}+1,&x<c,\\-1,&x\geq c,\end{cases}
 \qquad c=1.25892909961 .
\end{equation}
The implemented mixture has one numerically resolved Bayes boundary.
This is not a global theorem about arbitrary mixture cards.
At the nominal card $P_{\rm wrong}=.02202246905$; the majority-only
error is $.22530448517$. The chosen overall error limit $.05$ is an
engineering convention, not a statistical discovery threshold.
The nominal false-minus rate given $j=+1$ is $1.0315\%$; the
false-plus rate given $j=-1$ is $6.2279\%$. Passing the overall
criterion does not mean that both signs have at least 95\% recall.

The field $X$, previously called a detector, is a \emph{recoil partner};
its action does not yet specify an instrument that reads its state.
Likewise the G1 external sources describe linear response, not a
calibrated tracking or calorimeter system. A physical energy unit,
apparatus couplings, mass identification, acceptance and calibration
data are needed before claiming hardware feasibility. In particular,
the external model $U(1)$ has not been identified with electromagnetism.
The dependence of real momentum measurements on charge, magnetic field,
geometry and tracking errors is illustrated by \cite{pdg}; those
apparatus inputs are not implied by this prototype.

\section{Unknown calibration: keep the classifier frozen}
Model the true measured estimator as
\begin{equation}
 x=(1+a)u+b+\xi,\qquad
 u=p_{f,sq}(p),\qquad \xi\sim{\cal N}(0,\sigma^2).
\end{equation}
Here $a$ is a fractional scale error, $b$ an offset in model momentum
units, and $\sigma$ the actual noise width. Declare the stress domain
\begin{equation}\label{eq:box}
 |a|\leq .01,\qquad |b|\leq .02,\qquad .04\leq\sigma\leq .06 .
\end{equation}
These are test tolerances, not measured uncertainties. For the frozen
rule \eqref{eq:frozen}, the conditional error is exactly
\begin{align}
 R_0(a,b,\sigma)=\frac1Z\sum_q\int_0^\infty\dd p\, f_{.1}(p)
 \bigg[&
 w_0K_{0,0,q}(p)\,
 \Phi_{\!N}\!\left(\frac{(1+a)u_{+q}+b-c}{\sigma}\right)\nonumber\\
 &+w_{-2}K_{-2,0,q}(p)\,
 \Phi_{\!N}\!\left(\frac{c-(1+a)u_{-q}-b}{\sigma}\right)\bigg].
 \label{eq:risk}
\end{align}
The two terms separately give the joint probabilities for false minus
and false plus. Dividing by $\pi_+$ or $\pi_-$ gives the respective
conditional error rates; posterior purity divides by the predicted
class's total probability instead. Low aggregate error is not identical
to equal performance on both signs.

If $a,b$ were known, the invertible transformation
$y=(x-b)/(1+a)$ would reduce the oracle problem to a Gaussian noise
width $\sigma/(1+a)$:
\begin{equation}
 R_{\rm oracle}(a,b,\sigma)
 =R_{\rm Bayes}\!\left(\frac{\sigma}{1+a}\right)
 \leq R_0(a,b,\sigma).
\end{equation}
Thus the oracle is independent of known $b$. Reoptimizing the decision
rule separately at each unknown calibration value would hide the
actual robustness problem, so it is reported only as a lower bound.

For a rectangular sub-box $[a_-,a_+]\times[b_-,b_+]\times
[\sigma_-,\sigma_+]$, set
\begin{equation}
 z_+(u)=(1+a_+)u+b_+-c,\qquad
 z_-(u)=c-(1+a_-)u-b_- .
\end{equation}
Because $u\geq0$ and $\Phi_{\!N}$ is increasing, each error term is
bounded by
\begin{equation}
 B(z)=
 \begin{cases}
 \Phi_{\!N}(z/\sigma_-),&z\geq0,\\
 \Phi_{\!N}(z/\sigma_+),&z<0 .
 \end{cases}
\end{equation}
Replacing the two CDFs in \eqref{eq:risk} by these bounds gives a
conservative upper bound on the risk throughout that sub-box.
Different components are allowed different worst nuisances, so it
may be loose. A finite cover $\{C_k\}$ yields
$\sup_C R_0\leq\max_k U(C_k)$, not a sum of the cell bounds.
The numerical scan maximum is only a lower bound on this supremum.
CDF integrals are quadrature-checked; analytic inequalities do not
turn floating-point integration into an interval-arithmetic proof.

For this card the worst of the 27 declared grid points gives about
$3.8648\%$ at $(a,b,\sigma)=(-.01,-.02,.06)$. The unsplit bound is
$5.3902\%$: this loose bound alone cannot certify the 5\% target,
but exceeding the target is not an actual failure. One predetermined
$2\times2\times2$ cover of \eqref{eq:box} tightens the envelope to
about $4.46\%$, sufficient without repeated parameter fitting or
finer scans. The implementation splits the incident quadrature at
CDF-envelope kinks and checks the integral separately.

\section{Alternative: a noisy independent sideband tag}
Retain $\sigma_i=.1$ and compare the coarser $\sigma_d=.1$ readout
with a hypothetical extra label $z\in\{-1,0,1\}$. Declare
\begin{equation}
 T_\eta(z\mid q)=(1-\eta)\delta_{zq}+\frac{\eta}{3},\qquad
 0\leq\eta\leq1,\qquad P(z\ne q)=\frac{2\eta}{3}.
\end{equation}
The parameter $\eta$ is mixing strength, not the mistag probability.
This assumes the tag is independent of momentum noise and incoming
spread conditional on true $q$. It is a candidate response contract,
not a measurement derived from the prescribed classical pump.
The joint likelihood and optimal error are
\begin{equation}
 H_s(x,z)=\sum_qT_\eta(z\mid q)h_{sq}(x),\qquad
 e_\eta=\frac1Z\sum_z\int\dd x\,\min(H_+(x,z),H_-(x,z)).
\end{equation}
For $\eta_2\geq\eta_1$, $\eta_1<1$, define
$\lambda=(\eta_2-\eta_1)/(1-\eta_1)$. The uniform matrix
$U_{zq}=1/3$ obeys $U^2=U$, giving
\begin{equation}
 T_{\eta_2}=T_\lambda T_{\eta_1}.
\end{equation}
Additional garbling cannot improve optimal inference:
$\sum_z|\sum_{z'}T_\lambda(z\mid z')d_{z'}|
\leq\sum_{z'}|d_{z'}|$. Apply this pointwise to
$d_z=H_+(x,z)-H_-(x,z)$ to prove
\begin{equation}
e_{\rm ideal}=e_0\leq e_{\eta_1}\leq e_{\eta_2}
\leq e_1=e_{\rm no\ tag}.
\end{equation}
The code verifies both limiting cases and intermediate noise levels.
At $\sigma_d=.1$, the passing-set endpoint lies in
$\eta\in[.1758,.1768]$, corresponding to wrong-tag probability
$11.72$--$11.78\%$ \emph{within this specific response family}.
The endpoints bracket a numerical crossing, not calibration uncertainty.
At the primary $\sigma_d=.05$, no tag is required for the overall
5\% criterion.

\paragraph{Why aggregate tag accuracy is insufficient.}
In this selected ensemble $P(q=0)=.88983936$. A useless tag that
always emits $z=0$ therefore has $88.98\%$ raw accuracy, yet leaves
the coarser-momentum error at $5.94485\%$. This exceeds the roughly
88.2\% correct-label fraction at the symmetric-channel threshold.
There is no contradiction: the constant tag has a different
confusion matrix and no information. The requirement is the full
conditional response $T(z\mid q)$, not a single accuracy percentage.
These tag comparisons assume a known calibrated response and its
optimal classifier; they are not a frozen-rule robustness guarantee.

\subsection*{Do not count the same information twice}
Suppose the proposed tag is merely a deterministic or noisy processing
of the same measured $x$: $H_s(x,z)=h_s(x)T(z\mid x)$.
Then, identically,
\begin{equation}
 \sum_z\min(h_+(x)T(z\mid x),h_-(x)T(z\mid x))
 =\min(h_+(x),h_-(x)).
\end{equation}
The Bayes error cannot improve. Computing a pump label from the same
outgoing momentum and calling it an independent tag is therefore
invalid. A drive clock records phase, not the event's absorbed energy.
Even a genuinely additional sensor needs its joint response and
correlations modeled; it does not automatically satisfy $T_\eta$.
For two distinct sensors sharing the same true momentum, the correct
likelihood integrates that latent momentum once:
$h_s(x,y)=\sum_qw_{l_s}\int fK_{l_s,0,q}R_1(x\mid u)R_2(y\mid u)\dd p$.
Multiplying separately marginalized $h_s(x)h_s(y)$ is generally wrong.

\section{Closure and experimental handoff}
The accompanying reproducible reports give the nominal confusion
matrix, frozen calibration risks, conservative box bounds and noisy-tag
requirements. No action parameters are retuned. The primary choice
remains the momentum-only contract; the noisy tag is an alternative
requirement if finer momentum resolution is unavailable.

The bounded mathematical step ends here. To proceed physically, choose
the energy scale and readout interactions or supply an actual system's
response/beam calibration. Do not replace this missing physical choice
with indefinitely finer Gaussian scans. Absolute collision probability
is still unnecessary for this conditional discrimination question.
Only a request for absolute yields, finite-pulse effects or coherent
interference calls for four-dimensional packets, timing and a pump
envelope.

\begin{thebibliography}{1}
\bibitem{pdg} Particle Data Group, \emph{Particle Detectors at
Accelerators}, Review of Particle Physics (2024), section 35.13,
\href{https://pdg.lbl.gov/2024/reviews/rpp2024-rev-particle-detectors-accel.pdf}
{Official review (PDF)}.
\end{thebibliography}
\end{document}
